% paper/sec00-example.tex -- FORMAT EXAMPLE, not a section of any book.
%
% Part II reads the paper section by section (CLAUDE.md, Phase 5).  Copy this
% file to paper/secNN-SLUG.tex.  Each file is a \chapter: Part II chapters
% are numbered after Part I, which is why the running head is set by hand.
% main.tex deliberately does NOT \include it.
%
% ASCII only under tex/.
\chapter{<Title: what this section of the paper achieves>}
\chaptermark{<Short running head>}
\label{ch:read-SLUG}

% Same four-item header as a Part I chapter, read differently: here "Assumed"
% points at the preparatory chapters, and "How far things are proved" says
% which of the paper's arguments are read as printed and which had to be
% repaired.
\begin{chapterheader}
\textbf{Purpose.}
To read \S<n> of the paper. <What comes out of the section and nothing else
does -- name the two or three objects or inequalities that survive it.>

\medskip
\textbf{Assumed.}
Chapter~\ref{ch:SLUG} for <what>; <further preparatory chapters, each with
what it supplies>. <Anything assumed that Part I did not cover, named here.>

\medskip
\textbf{Supports.}
\S<n> of the paper, pp.~<a>--<b>. <Where its output is consumed: which later
section, and for what.>

\medskip
\textbf{How far things are proved.}
<The default is that the paper's proofs are read rather than reproduced.  Name
the exceptions: the places where the printed argument does not close, what was
done about each, and that the repair is marked in the text.  If a statement of
the paper is wrong rather than compressed, say so here and record it in the
list of misprints.>
\end{chapterheader}

%% =====================================================================
\section{<What the section replaces, or the problem it opens with>}
\label{sec:read-SLUG-opening}

<Set the scene: what the paper has in hand at the start of this section, and
what it cannot yet do.  A section of Part II opens by saying why the section
exists, not by summarising it.>

%% =====================================================================
\section{<The argument, in the paper's order>}
\label{sec:read-SLUG-argument}

<Follow the paper.  Quote its numbering as text -- Lemma <n.m>, equation
(<n.m>) -- so that the reader can put the two documents side by side.  Those
numbers pass through no cross-reference and nothing checks them
automatically, which is why
the project keeps a source-number checker (tex/<paper>-check.sh).>

\begin{lem}[{\cite[Lem.~<n.m>]{Har77}}]
\label{lem:read-SLUG}
<Restate a statement of the paper only when the reading needs it in front of
the reader.  Otherwise cite it.>
\end{lem}

<Where the paper compresses a step, restore it and mark the restoration:>
\OWNPROOF{the paper obtains this in one phrase; the two lines below are ours}

<Where our own routine verification is added for the reader's convenience,
mark it more quietly:>
\OWNCHECK{a direct computation, no gap in the source}

<Where something remains genuinely unclear, say so rather than hedging:>
\uncertain{<what exactly is unclear, and what would settle it>}

\begin{rem}[what \S<n> hands to \S<n+1>]
\label{rem:read-SLUG-handover}
<The last remark of a Part II chapter states the situation at the end of the
section: the objects now available, the inequalities now known, and -- this is
the useful half -- what is NOT used again.>
\end{rem}

%% =====================================================================
\section{Summary: where this section is used}
\label{sec:read-SLUG-summary}

\begin{center}
\small
\begin{tabular}{@{}L{4.6cm}lL{5.9cm}@{}}
\toprule
Item & Here & Used in <PAPER> \\
\midrule
<the object the section constructs>
  & \ref{lem:read-SLUG}
  & <\S of the paper that consumes it, and for what> \\
<the restored argument>
  & \ref{rem:read-SLUG-handover}
  & <where the restoration matters later, if it does> \\
\bottomrule
\end{tabular}
\end{center}
